% ============================================================
\chapter{Rappels de Probabilit\'es}
\label{ch:rappels}
% ============================================================

Avant de plonger dans l'\'etude des processus stochastiques --- ces
fonctions al\'eatoires qui \'evoluent dans le temps --- il est indispensable
de s'assurer que nos fondations probabilistes sont solides. Ce chapitre
rappelle les outils essentiels~: espaces de probabilit\'e, variables
al\'eatoires, esp\'erance conditionnelle et convergences. Si le lecteur
ma\^itrise d\'ej\`a ces notions, il peut parcourir ce chapitre rapidement en
guise de r\'ef\'erence~; sinon, c'est ici que tout commence.

% ------------------------------------------------------------
\section{Espaces de probabilit\'e}
% ------------------------------------------------------------

Rappelons les trois ingr\'edients fondamentaux de la th\'eorie de Kolmogorov.

\begin{definition}[Tribu]
Soit $\Omega$ un ensemble non vide. Une \emph{tribu} (ou $\sigma$-alg\`ebre)
$\mathcal{F}$ sur $\Omega$ est une collection de sous-ensembles de $\Omega$
v\'erifiant~:
\begin{enumerate}
  \item $\Omega \in \mathcal{F}$,
  \item si $A \in \mathcal{F}$, alors $A^c \in \mathcal{F}$ (stabilit\'e par compl\'ementation),
  \item si $(A_n)_{n \geq 1} \subset \mathcal{F}$, alors $\bigcup_{n=1}^{\infty} A_n \in \mathcal{F}$
    (stabilit\'e par union d\'enombrable).
\end{enumerate}
Le couple $(\Omega, \mathcal{F})$ est appel\'e \emph{espace mesurable}.
\end{definition}

\begin{definition}[Espace de probabilit\'e]
Un \emph{espace de probabilit\'e} est un triplet $(\Omega, \mathcal{F}, \PP)$ o\`u~:
\begin{itemize}
  \item $\Omega$ est l'\emph{univers} (ensemble de tous les r\'esultats possibles),
  \item $\mathcal{F}$ est une tribu sur $\Omega$ (les \'ev\'enements),
  \item $\PP : \mathcal{F} \to [0,1]$ est une \emph{mesure de probabilit\'e}, c'est-\`a-dire
    une mesure $\sigma$-additive avec $\PP(\Omega) = 1$.
\end{itemize}
\end{definition}

\begin{remarque}
La $\sigma$-additivit\'e de $\PP$ signifie~: pour toute suite $(A_n)_{n \geq 1}$
d'\'ev\'enements deux \`a deux disjoints,
\[
  \PP\Bigl(\bigcup_{n=1}^{\infty} A_n\Bigr) = \sum_{n=1}^{\infty} \PP(A_n).
\]
\end{remarque}

\begin{definition}[Tribu engendr\'ee]
Soit $\mathcal{C}$ une collection de sous-ensembles de $\Omega$. La \emph{tribu
engendr\'ee} par $\mathcal{C}$, not\'ee $\sigma(\mathcal{C})$, est la plus petite
tribu contenant $\mathcal{C}$~:
\[
  \sigma(\mathcal{C}) = \bigcap \bigl\{ \mathcal{F} : \mathcal{F} \text{ tribu sur }
  \Omega,\; \mathcal{C} \subset \mathcal{F} \bigr\}.
\]
En particulier, la \emph{tribu bor\'elienne} $\mathcal{B}(\R)$ est $\sigma$-alg\`ebre
engendr\'ee par les ouverts de $\R$.
\end{definition}

\section{Variables al\'eatoires}

\begin{definition}[Variable al\'eatoire]
Soit $(\Omega, \mathcal{F}, \PP)$ un espace de probabilit\'e. Une \emph{variable
al\'eatoire} (v.a.) est une application mesurable $X : (\Omega, \mathcal{F}) \to
(\R, \mathcal{B}(\R))$, c'est-\`a-dire~:
\[
  \forall B \in \mathcal{B}(\R), \quad X^{-1}(B) = \{\omega \in \Omega : X(\omega) \in B\} \in \mathcal{F}.
\]
Plus g\'en\'eralement, un \emph{vecteur al\'eatoire} est une application mesurable
$X : (\Omega, \mathcal{F}) \to (\R^d, \mathcal{B}(\R^d))$.
\end{definition}

\begin{definition}[Tribu engendr\'ee par une v.a.]
La tribu engendr\'ee par une variable al\'eatoire $X$ est
\[
  \sigma(X) = \{X^{-1}(B) : B \in \mathcal{B}(\R)\}.
\]
C'est la plus petite tribu rendant $X$ mesurable. Intuitivement, $\sigma(X)$
repr\'esente l'\emph{information} contenue dans l'observation de $X$.
\end{definition}

\begin{theoreme}[Lemme de Doob]
\label{thm:doob_lemme}
Soit $X$ une variable al\'eatoire et $\mathcal{G} = \sigma(X)$. Une variable
al\'eatoire $Y$ est $\mathcal{G}$-mesurable si et seulement s'il existe une
fonction bor\'elienne $g : \R \to \R$ telle que $Y = g(X)$.
\end{theoreme}

\section{Esp\'erance et int\'egration}

\begin{definition}[Esp\'erance]
L'\emph{esp\'erance} d'une variable al\'eatoire $X \geq 0$ est
\[
  \E[X] = \int_{\Omega} X(\omega) \, d\PP(\omega).
\]
Pour une v.a.\ g\'en\'erale, $\E[X] = \E[X^+] - \E[X^-]$ lorsque l'une au
moins des deux quantit\'es est finie, o\`u $X^+ = \max(X, 0)$ et $X^- = \max(-X, 0)$.
On dit que $X$ est \emph{int\'egrable} si $\E[\abs{X}] < \infty$.
\end{definition}

\begin{theoreme}[Convergence domin\'ee de Lebesgue]
\label{thm:conv_dom}
Soit $(X_n)_{n \geq 1}$ une suite de v.a.\ convergeant p.s.\ vers $X$.
S'il existe une v.a.\ int\'egrable $Y$ telle que $\abs{X_n} \leq Y$ p.s.\ pour tout $n$,
alors $X$ est int\'egrable et
\[
  \lim_{n \to \infty} \E[X_n] = \E[X].
\]
\end{theoreme}

\begin{theoreme}[Convergence monotone]
\label{thm:conv_mon}
Soit $(X_n)_{n \geq 1}$ une suite croissante de v.a.\ positives~:
$0 \leq X_1 \leq X_2 \leq \cdots$ p.s. Alors
\[
  \lim_{n \to \infty} \E[X_n] = \E\bigl[\lim_{n \to \infty} X_n\bigr].
\]
\end{theoreme}

\begin{theoreme}[Lemme de Fatou]
Soit $(X_n)_{n \geq 1}$ une suite de v.a.\ positives. Alors
\[
  \E\bigl[\liminf_{n \to \infty} X_n\bigr] \leq \liminf_{n \to \infty} \E[X_n].
\]
\end{theoreme}

\section{Espaces $L^p$ et in\'egalit\'es}

\begin{definition}[Espace $L^p$]
Pour $p \in [1, \infty)$, l'espace $L^p(\Omega, \mathcal{F}, \PP)$ est l'ensemble
des (classes de) v.a.\ $X$ telles que $\E[\abs{X}^p] < \infty$, muni de la norme
\[
  \norm{X}_{L^p} = \bigl(\E[\abs{X}^p]\bigr)^{1/p}.
\]
L'espace $L^{\infty}$ est l'ensemble des v.a.\ essentiellement born\'ees.
\end{definition}

\begin{theoreme}[In\'egalit\'e de H\"older]
Soient $p, q \in (1, \infty)$ avec $\frac{1}{p} + \frac{1}{q} = 1$. Pour
$X \in L^p$ et $Y \in L^q$,
\[
  \E[\abs{XY}] \leq \norm{X}_{L^p} \norm{Y}_{L^q}.
\]
Le cas $p = q = 2$ donne l'\emph{in\'egalit\'e de Cauchy--Schwarz}.
\end{theoreme}

\begin{theoreme}[In\'egalit\'e de Jensen]
\label{thm:jensen}
Si $X$ est int\'egrable et $\varphi : \R \to \R$ est convexe, alors
\[
  \varphi(\E[X]) \leq \E[\varphi(X)],
\]
pourvu que $\E[\abs{\varphi(X)}] < \infty$.
\end{theoreme}

\section{Esp\'erance conditionnelle}

\begin{definition}[Esp\'erance conditionnelle]
\label{def:esp_cond}
Soit $X \in L^1(\Omega, \mathcal{F}, \PP)$ et $\mathcal{G} \subset \mathcal{F}$
une sous-tribu. L'\emph{esp\'erance conditionnelle} de $X$ sachant $\mathcal{G}$,
not\'ee $\E[X \mid \mathcal{G}]$, est l'unique (p.s.) variable al\'eatoire v\'erifiant~:
\begin{enumerate}
  \item $\E[X \mid \mathcal{G}]$ est $\mathcal{G}$-mesurable,
  \item pour tout $G \in \mathcal{G}$, $\displaystyle \int_G \E[X \mid \mathcal{G}] \, d\PP
    = \int_G X \, d\PP$.
\end{enumerate}
L'existence est garantie par le th\'eor\`eme de Radon--Nikodym.
\end{definition}

\begin{theoreme}[Propri\'et\'es de l'esp\'erance conditionnelle]
\label{thm:prop_ec}
Soient $X, Y \in L^1$ et $\mathcal{G}, \mathcal{H}$ des sous-tribus de $\mathcal{F}$.
\begin{enumerate}
  \item \textbf{Lin\'earit\'e}~: $\E[aX + bY \mid \mathcal{G}] = a\E[X \mid \mathcal{G}]
    + b\E[Y \mid \mathcal{G}]$ p.s.
  \item \textbf{Positivit\'e}~: si $X \geq 0$ p.s., alors $\E[X \mid \mathcal{G}] \geq 0$ p.s.
  \item \textbf{Tour de Babel (lissage)}~: si $\mathcal{H} \subset \mathcal{G}$,
    alors $\E[\E[X \mid \mathcal{G}] \mid \mathcal{H}] = \E[X \mid \mathcal{H}]$ p.s.
  \item \textbf{Factorisation}~: si $Y$ est $\mathcal{G}$-mesurable et born\'ee,
    $\E[XY \mid \mathcal{G}] = Y \E[X \mid \mathcal{G}]$ p.s.
  \item \textbf{Ind\'ependance}~: si $\sigma(X)$ et $\mathcal{G}$ sont ind\'ependantes,
    $\E[X \mid \mathcal{G}] = \E[X]$ p.s.
  \item \textbf{Jensen conditionnel}~: si $\varphi$ est convexe et $\varphi(X) \in L^1$,
    $\varphi(\E[X \mid \mathcal{G}]) \leq \E[\varphi(X) \mid \mathcal{G}]$ p.s.
\end{enumerate}
\end{theoreme}

\begin{intuition}[title=\textbf{Interpr\'etation g\'eom\'etrique}]
Dans $L^2(\Omega, \mathcal{F}, \PP)$, l'esp\'erance conditionnelle $\E[X \mid \mathcal{G}]$
est la \emph{projection orthogonale} de $X$ sur le sous-espace ferm\'e
$L^2(\Omega, \mathcal{G}, \PP)$. Cela justifie la propri\'et\'e~:
\[
  \E\bigl[\bigl(X - \E[X \mid \mathcal{G}]\bigr)^2\bigr]
  \leq \E\bigl[(X - Z)^2\bigr]
\]
pour tout $Z \in L^2(\Omega, \mathcal{G}, \PP)$.
\end{intuition}

\section{Ind\'ependance}

\begin{definition}[Ind\'ependance de tribus]
Des sous-tribus $\mathcal{G}_1, \ldots, \mathcal{G}_n$ de $\mathcal{F}$ sont
\emph{ind\'ependantes} si pour tout choix $A_i \in \mathcal{G}_i$~:
\[
  \PP(A_1 \cap \cdots \cap A_n) = \PP(A_1) \cdots \PP(A_n).
\]
Des v.a.\ $X_1, \ldots, X_n$ sont ind\'ependantes si les tribus $\sigma(X_1), \ldots,
\sigma(X_n)$ sont ind\'ependantes.
\end{definition}

\begin{definition}[Ind\'ependance mutuelle]
Une famille $(X_i)_{i \in I}$ de v.a.\ est \emph{mutuellement ind\'ependante}
si pour tout sous-ensemble fini $J \subset I$, les v.a.\ $(X_j)_{j \in J}$
sont ind\'ependantes.
\end{definition}

\begin{theoreme}[Loi $0$--$1$ de Kolmogorov]
\label{thm:kolmo01}
Soit $(X_n)_{n \geq 1}$ une suite de v.a.\ ind\'ependantes. La \emph{tribu
terminale} (ou queue)
\[
  \mathcal{T} = \bigcap_{n=1}^{\infty} \sigma(X_n, X_{n+1}, \ldots)
\]
est \emph{triviale}~: pour tout $A \in \mathcal{T}$, $\PP(A) \in \{0, 1\}$.
\end{theoreme}

\section{Modes de convergence}

\begin{definition}[Convergences]
Soit $(X_n)_{n \geq 1}$ une suite de v.a.\ et $X$ une v.a.\ d\'efinies sur
$(\Omega, \mathcal{F}, \PP)$.
\begin{enumerate}
  \item \textbf{Convergence presque s\^ure}~: $X_n \xrightarrow{\text{p.s.}} X$
    si $\PP(\{\omega : X_n(\omega) \to X(\omega)\}) = 1$.
  \item \textbf{Convergence dans $L^p$}~: $X_n \xrightarrow{L^p} X$ si
    $\E[\abs{X_n - X}^p] \to 0$.
  \item \textbf{Convergence en probabilit\'e}~: $X_n \xrightarrow{\PP} X$ si
    pour tout $\varepsilon > 0$, $\PP(\abs{X_n - X} > \varepsilon) \to 0$.
  \item \textbf{Convergence en loi}~: $X_n \xrightarrow{\mathcal{L}} X$ si
    $\E[f(X_n)] \to \E[f(X)]$ pour toute $f : \R \to \R$ continue born\'ee.
\end{enumerate}
\end{definition}

\begin{formules}[title=\textbf{Hi\'erarchie des convergences}]
\[
  \text{p.s.} \implies \text{en probabilit\'e} \implies \text{en loi}
\]
\[
  L^p \implies L^q \;\;(q \leq p) \implies \text{en probabilit\'e} \implies \text{en loi}
\]
Les r\'eciproques sont fausses en g\'en\'eral. Cependant~:
\begin{itemize}
  \item $X_n \xrightarrow{\PP} X$ $\implies$ il existe une sous-suite $X_{n_k} \xrightarrow{\text{p.s.}} X$.
  \item Si $X_n \xrightarrow{\mathcal{L}} c$ (constante), alors $X_n \xrightarrow{\PP} c$.
\end{itemize}
\end{formules}

\section{Lois des grands nombres}

\begin{theoreme}[Loi forte des grands nombres]
\label{thm:lfgn}
Soit $(X_n)_{n \geq 1}$ une suite de v.a.\ i.i.d.\ int\'egrables, de moyenne
$\mu = \E[X_1]$. Alors
\[
  \frac{1}{n} \sum_{k=1}^{n} X_k \xrightarrow{\text{p.s.}} \mu.
\]
\end{theoreme}

\begin{theoreme}[Th\'eor\`eme central limite]
\label{thm:tcl}
Soit $(X_n)_{n \geq 1}$ une suite de v.a.\ i.i.d.\ avec $\E[X_1] = \mu$ et
$\Var(X_1) = \sigma^2 \in (0, \infty)$. Alors
\[
  \frac{1}{\sqrt{n}} \sum_{k=1}^{n} \frac{X_k - \mu}{\sigma}
  \xrightarrow{\mathcal{L}} \mathcal{N}(0, 1).
\]
\end{theoreme}

\section{Lemme de Borel--Cantelli}

\begin{theoreme}[Lemme de Borel--Cantelli]
\label{thm:borel_cantelli}
Soit $(A_n)_{n \geq 1}$ une suite d'\'ev\'enements.
\begin{enumerate}
  \item Si $\sum_{n=1}^{\infty} \PP(A_n) < \infty$, alors
    $\PP(\limsup_{n} A_n) = 0$, c'est-\`a-dire p.s.\ un nombre fini de $A_n$
    se r\'ealisent.
  \item Si les $A_n$ sont \emph{ind\'ependants} et $\sum_{n=1}^{\infty} \PP(A_n) = \infty$,
    alors $\PP(\limsup_{n} A_n) = 1$, c'est-\`a-dire p.s.\ une infinit\'e de $A_n$
    se r\'ealisent.
\end{enumerate}
\end{theoreme}

\begin{attention}[title=\textbf{Hypoth\`ese d'ind\'ependance}]
La seconde partie du lemme de Borel--Cantelli n\'ecessite l'ind\'ependance des
\'ev\'enements. Sans cette hypoth\`ese, la divergence de la s\'erie $\sum \PP(A_n)$
ne suffit pas \`a garantir $\PP(\limsup A_n) = 1$.
\end{attention}

\section{Fonctions caract\'eristiques}

\begin{definition}[Fonction caract\'eristique]
La \emph{fonction caract\'eristique} d'une v.a.\ $X$ est
\[
  \varphi_X(t) = \E[e^{itX}], \quad t \in \R.
\]
\end{definition}

\begin{proposition}[Propri\'et\'es]
\begin{enumerate}
  \item $\varphi_X(0) = 1$ et $\abs{\varphi_X(t)} \leq 1$ pour tout $t$.
  \item $\varphi_X$ est uniform\'ement continue sur $\R$.
  \item $\varphi_X$ caract\'erise la loi de $X$~: si $\varphi_X = \varphi_Y$,
    alors $X \overset{\mathcal{L}}{=} Y$.
  \item Si $X$ et $Y$ sont ind\'ependantes, $\varphi_{X+Y}(t) = \varphi_X(t) \varphi_Y(t)$.
  \item Si $\E[\abs{X}^n] < \infty$, $\varphi_X^{(n)}(0) = i^n \E[X^n]$.
\end{enumerate}
\end{proposition}

\begin{theoreme}[Th\'eor\`eme de continuit\'e de L\'evy]
\label{thm:levy}
$X_n \xrightarrow{\mathcal{L}} X$ si et seulement si $\varphi_{X_n}(t) \to \varphi_X(t)$
pour tout $t \in \R$.
\end{theoreme}

\section{Exercices}

\begin{exercice}
Soit $(\Omega, \mathcal{F}, \PP)$ un espace de probabilit\'e et $X$ une v.a.\ positive.
Montrer que $\E[X] = \int_0^{\infty} \PP(X > t) \, dt$.
\end{exercice}

\begin{exercice}
Montrer que la convergence p.s.\ n'implique pas la convergence $L^1$ en construisant
un contre-exemple.
\end{exercice}

\begin{exercice}
Soit $(X_n)_{n \geq 1}$ une suite de v.a.\ ind\'ependantes avec
$\PP(X_n = n) = 1/n$ et $\PP(X_n = 0) = 1 - 1/n$. Montrer que $X_n \to 0$ p.s.
\end{exercice}

\begin{exercice}
D\'emontrer la propri\'et\'e de lissage de l'esp\'erance conditionnelle
(th\'eor\`eme~\ref{thm:prop_ec}, point~3).
\end{exercice}

\begin{exercice}
Soit $X$ une v.a.\ int\'egrable et $\mathcal{G}_n \uparrow \mathcal{G}_{\infty}$
une filtration croissante de sous-tribus. Montrer que
$\E[X \mid \mathcal{G}_n] \to \E[X \mid \mathcal{G}_{\infty}]$ p.s.\ et dans $L^1$.
\emph{(C'est le th\'eor\`eme de convergence des martingales de L\'evy.)}
\end{exercice}
